\chapter*{LỜI NÓI ĐẦU}
\addcontentsline{toc}{chapter}{LỜI NÓI ĐẦU}

Trong kỷ nguyên số hóa hiện nay, việc ứng dụng công nghệ thông tin vào quản lý thư viện không còn là xu hướng mà đã trở thành nhu cầu thiết yếu đối với mọi cơ sở giáo dục và văn hóa. Một hệ thống quản lý thư viện hiện đại cần đáp ứng được các yêu cầu về tốc độ xử lý, khả năng mở rộng, tính bảo mật và trải nghiệm người dùng trên nhiều nền tảng khác nhau. Xuất phát từ thực tiễn đó, nhóm em đã lựa chọn đề tài “Xây dựng hệ thống quản lý thư viện” theo kiến trúc microservices làm bài tập lớn cho học phần Thiết kế \& Kiến trúc phần mềm. Hệ thống được phân tách thành ba dịch vụ độc lập: Dịch vụ Danh mục (Catalog Service – N1) phụ trách quản lý kho sách và thông tin đầu sách, Dịch vụ Lưu thông (Circulation Service – N2) đảm nhiệm toàn bộ nghiệp vụ mượn – trả – gia hạn – phí phạt và Dịch vụ Định danh (Identity Service – N3) chịu trách nhiệm xác thực và phân quyền người dùng. Trong đó, nhóm em được phân công phát triển Dịch vụ Lưu thông (N2) – thành phần trung tâm kết nối hai dịch vụ còn lại và xử lý toàn bộ luồng nghiệp vụ cốt lõi của thư viện. Báo cáo này trình bày chi tiết quá trình nhóm em phân tích, thiết kế và triển khai Service Lưu thông theo kiến trúc REST API trên nền tảng ASP.NET Core 8, kết hợp giao diện người dùng xây dựng bằng Vue 3. Báo cáo được cấu trúc thành ba chương: Chương 1 giới thiệu các khái niệm nền tảng về kiến trúc phần mềm và các kiểu kiến trúc phổ biến trong phát triển ứng dụng hiện đại, Chương 2 mô tả tổng quan bài toán quản lý thư viện và làm rõ phạm vi công việc của nhóm N2, Chương 3 đi sâu vào chi tiết triển khai mã nguồn và các thử nghiệm thay đổi tham số hệ thống. Trong quá trình thực hiện, nhóm em đã vận dụng các nguyên lý thiết kế phần mềm như nguyên lý đơn nhiệm (Single Responsibility), nguyên lý đóng mở (Open/Closed) và các mẫu kiến trúc như Dependency Injection, Repository Pattern cùng Event-Driven Architecture để xây dựng một hệ thống có tính module hóa cao, dễ bảo trì và sẵn sàng mở rộng khi cần tích hợp thêm dịch vụ mới. Nhóm em xin gửi lời cảm ơn chân thành đến thầy ThS. Đỗ Quang Thơ – giảng viên học phần Thiết kế \& Kiến trúc phần mềm – đã tận tình hướng dẫn và cung cấp những kiến thức nền tảng quý báu để nhóm có thể hoàn thành bài tập lớn này. Mặc dù đã cố gắng
